\section{Conclusion}
This paper develops a vision- and tension-informed connection between inertial identification and supporting formation geometry for cooperative aerial transport through known obstacles. The payload pose reference family is prescribed; conic optimization selects the cable forces, compliant support endpoints and a feasible timing candidate. A reconstructed surface anchors markerless pose measurements and a nonnegative representation of inertial moments. Integral momentum observations identify the CoM and inertia used by a conic cable-force planner. Compliant endpoint construction and sequential separation then connect these forces to the aircraft formation throughout rotation and aperture passage.

Three-aircraft and four-aircraft rectangular/irregular configurations complete the base task with 2.62--4.10 cm position errors and $0.66$--$2.64^\circ$ attitude errors. The mass-matched five-arm comparison attributes the main base-task tracking gain to CoM correction, while complete-inertia accuracy and transport-only prediction quantify the additional identified moments. A forked payload subsequently completes a 28 m course with successive side-on, horizontal and opposite-side orientations, reaching approximately $\pm90^\circ$ roll while maintaining positive cable tensions and clearance. These findings support the use of identified internal mass moments when coordinating payload attitude and supporting formation through spatial restrictions.

\section*{Funding}
This work was supported in part by the Foundation for Innovative Research Groups of the National Natural Science Foundation of China under Grant 62421004, and in part by the National Natural Science Foundation of China under Grants 62303348 and 62173242.

\section*{Data availability}
The simulation and identification source code, task configurations, true and fitted inertial parameters, archived trajectories, ablation and repeated-seed experiments and figure-generation scripts are available at \url{https://github.com/shepherd-p/ForestSling}. The repository documents the analysis and rerun commands and identifies each experiment by its configuration and random seed. Recorded trajectory visualizations are provided with the reproducibility materials. Proprietary Unreal environment assets are excluded; numerical analysis and MuJoCo simulation do not require those assets.
